% source: RKF/theorum/thermodynamics/03_canonical_response_cost_barrier.md
\hypertarget{canonical-thermodynamic-response-cost-and-barrier-theorem}{%
\chapter{Canonical Thermodynamic Response-Cost and Barrier Theorem}\label{canonical-thermodynamic-response-cost-and-barrier-theorem}}

\hypertarget{canonical-response-frame}{%
\section{1. Canonical response frame}\label{canonical-response-frame}}

Let the four thermodynamic response columns be

\[
\rho_X(x)\in\mathcal Y,
\qquad X\in\{p,v,s,t\},
\]

with synthesis map

\[
T_x:\mathbb C^4\to\mathcal Y,
\qquad T_xe_X=\rho_X(x).
\]

Define the compatibility Gram matrix

\[
M_x=T_x^*T_x.
\]

On

\[
\mathcal S_{\mathrm{th}}(x)=(\ker M_x)^\perp,
\]

define

\[
\boxed{
J_{\mathrm{th}}(x)=T_xM_x^{-1/2}.
}
\]

Then

\[
J_{\mathrm{th}}(x)^*J_{\mathrm{th}}(x)=I_{\mathcal S_{\mathrm{th}}(x)}.
\]

\hypertarget{canonical-source-to-response-coupling}{%
\section{2. Canonical source-to-response coupling}\label{canonical-source-to-response-coupling}}

Let \(J_q\) be an isometric reconstruction of a declared native source block and let \(\mathcal C\) be the native channel isometry. Define

\[
\boxed{
B_q(x)=J_{\mathrm{th}}(x)^*\mathcal C J_q.
}
\]

Let

\[
H_{\mathrm{th}}(x)>0
\]

be a stable response metric obtained from positive susceptibilities or compatibility-Gram whitening.

Define the response cost

\[
\boxed{
Q_q(x)=B_q(x)^*H_{\mathrm{th}}(x)B_q(x)\ge0.
}
\]

\hypertarget{theorem-2.1-bounded-thermodynamic-response-barrier}{%
\section{Theorem 2.1 (bounded thermodynamic response barrier)}\label{theorem-2.1-bounded-thermodynamic-response-barrier}}

Define

\[
A_q=e^{-Q_q/2},
\qquad
\boxed{
W_q=I-e^{-Q_q}.
}
\]

Then

\[
\boxed{0\le W_q<I}
\]

and

\[
\boxed{
\ker W_q=\ker Q_q=\ker B_q.
}
\]

\hypertarget{proof}{%
\subsection{Proof}\label{proof}}

Since \(H_{\mathrm{th}}>0\),

\[
Q_q=(H_{\mathrm{th}}^{1/2}B_q)^*(H_{\mathrm{th}}^{1/2}B_q)\ge0,
\]

and

\[
\ker Q_q=\ker(H_{\mathrm{th}}^{1/2}B_q)=\ker B_q.
\]

Functional calculus on the nonnegative operator \(Q_q\) gives

\[
0<e^{-Q_q}\le I,
\]

hence \(0\le W_q<I\). The scalar function \(1-e^{-t}\) vanishes exactly at \(t=0\), so

\[
\ker W_q=\ker Q_q.
\]

QED.

\hypertarget{basis-invariance}{%
\section{3. Basis invariance}\label{basis-invariance}}

Under unitary changes of source and response coordinates,

\[
B_q\mapsto V^*B_qU,
\qquad
H_{\mathrm{th}}\mapsto V^*H_{\mathrm{th}}V.
\]

Therefore

\[
Q_q\mapsto U^*Q_qU,
\qquad
W_q\mapsto U^*W_qU.
\]

Spectrum, rank, kernel, and positivity are invariant.

\hypertarget{barrier-interpretation}{%
\section{4. Barrier interpretation}\label{barrier-interpretation}}

The spectrum of \(Q_q\) measures response cost. The bounded defect \(W_q\) is the corresponding response barrier. Radial response magnitude changes the barrier strength; unitary angular changes rotate its orientation without changing its spectrum.

\hypertarget{claim-boundary}{%
\section{Claim boundary}\label{claim-boundary}}

\begin{verbatim}
CANONICAL RESPONSE-FRAME WHITENING         PROVED
CANONICAL COUPLING B_q                     PROVED ON DECLARED SOURCE RANGE
POSITIVE RESPONSE COST Q_q                 PROVED
BOUNDED BARRIER W_q                        PROVED
KERNEL IDENTITY                            PROVED
UNITARY BASIS INVARIANCE                   PROVED
PHYSICAL SOURCE IDENTIFICATION             OPEN
CONTINUUM RESPONSE-CURVATURE IDENTITY       OPEN
\end{verbatim}

\begin{flushright}\small\texttt{RKF/theorum/thermodynamics/03\_canonical\_response\_cost\_barrier.md}\end{flushright}
